\itemhead{习题3.1 · 第5题（教材p.77）}{XJ-B1-C03-HW-P077-Y0403}
\label{XJ-B1-C03-HW-P077-Y0403}
{\footnotesize 湘教版教材p.77 / PDF第84页 / 习题3.1 / 第5题}
\field{原题}{下列各组函数中，表示同一个函数的是（\quad）
\choice{A}$y=x-1$ 和 $y=\dfrac{x^2-1}{x+1}$
\choice{B}$y=x^0$ 和 $y=1$
\choice{C}$y=x^2$ 和 $y=(x+1)^2$
\choice{D}$y=\dfrac{(\sqrt x)^2}x$ 和 $y=\dfrac x{(\sqrt x)^2}$}
\field{参考答案}{D。两函数都在 $(0,+\infty)$ 上恒为1。B项 $x^0$ 还要求 $x\ne0$。}
{\small \field{关联知识}{F01 函数本质、构成要素与相同函数（教材帮p.77、78 / PDF86、87）；F02 定义域与原式限制（教材帮p.79、80、81 / PDF88、89、90）}
\field{方法与易错点}{逐一核对每个允许输入是否有且仅有一个合乎目标集合的输出。 易错：把值域直接当作目标集合；只看解析式外观；认为常函数不是函数。}
\field{答案依据}{答案册相关题组：PDF第20、21、22页（册内页19、20、21）；按当前题库原题逐项重算，不按旧版题号机械对套。}}
